\documentclass[conference]{IEEEtran}
\IEEEoverridecommandlockouts

% Mantener tipografía Times tanto con pdfLaTeX como con XeLaTeX/Tectonic.
\usepackage{iftex}
\ifPDFTeX
  \usepackage[T1]{fontenc}
  \usepackage[utf8]{inputenc}
\else
  \usepackage{fontspec}
  \setmainfont{texgyretermes}[
    Extension=.otf, UprightFont=*-regular, BoldFont=*-bold,
    ItalicFont=*-italic, BoldItalicFont=*-bolditalic]
  \setmonofont{texgyrecursor}[
    Extension=.otf, UprightFont=*-regular, BoldFont=*-bold,
    ItalicFont=*-italic, BoldItalicFont=*-bolditalic]
\fi
\usepackage[spanish,es-tabla]{babel}

\usepackage{cite}
\usepackage{amsmath,amssymb,amsfonts}
\usepackage{algorithmic}
\usepackage{graphicx}
\usepackage{textcomp}
\usepackage{xcolor}
\def\BibTeX{{\rm B\kern-.05em{\sc i\kern-.025em b}\kern-.08em
    T\kern-.1667em\lower.7ex\hbox{E}\kern-.125emX}}
\begin{document}

\title{Predicción de resultados del fútbol uruguayo mediante métodos de aprendizaje supervisado}

\author{
\IEEEauthorblockN{Thiago Rivas}
\IEEEauthorblockA{C.I. 5.475.680-3 \\
\textit{Facultad de Ingeniería} \\
\textit{Universidad de la República}\\
Montevideo, Uruguay \\
thiago.rivas0021@gmail.com}
\and
\IEEEauthorblockN{Juan Duarte}
\IEEEauthorblockA{C.I. 5.201.302-9 \\
\textit{Facultad de Ingeniería} \\
\textit{Universidad de la República}\\
Montevideo, Uruguay \\
juan.duarte@fing.edu.uy}
\and
\IEEEauthorblockN{Agustin Occhiuzzi}
\IEEEauthorblockA{C.I. 5.480.562-4 \\
\textit{Facultad de Ingeniería} \\
\textit{Universidad de la República}\\
Montevideo, Uruguay \\
agustin.occhiuzzi@fing.edu.uy}
}

\maketitle

% =========================================================
% PRESUPUESTO DE ESPACIO (formato IEEE, 2 columnas, 10pt)
% Referencia: ~500 palabras/columna de solo texto, ~1000 palabras/página.
% Límite duro: 5 páginas TOTALES, incluyendo 1 tabla y 3 figuras ya
% reservadas en el esqueleto (ocupan aprox. 1.5 páginas equivalentes
% entre las tres + la tabla). Presupuesto de TEXTO detallado por
% sección/subsección en cada % TODO más abajo.
%
% RESUMEN POR SECCIÓN (para chequeo rápido de avance):
%   Abstract .......................  100-150 palabras
%   Introducción ...................  250-300 palabras
%   Trabajo relacionado (opcional) .   80-120 palabras (o se omite)
%   Dataset y preprocesamiento .....  550-650 palabras (6 subsecciones)
%   Modelos .........................  400-450 palabras (+ 3 ecuaciones)
%   Experimentación .................  180-220 palabras (+ 2 figuras)
%   Resultados ......................  350-400 palabras (+ tabla + fig.)
%   Conclusiones ....................  150-180 palabras
%   Uso de IA .......................   50-70 palabras
%   -------------------------------------------------------
%   TOTAL TEXTO ..................... ~2100-2600 palabras
%   + 1-2 tablas, 3-4 figuras -> ocupan el resto hasta las 5 páginas.
%
% Si al ir escribiendo una sección se pasa del número, recortar AHÍ
% antes de seguir con la siguiente — es más fácil ajustar sobre la
% marcha que recortar 500 palabras de golpe al final.
% =========================================================

\begin{abstract}
En este informe, se estudia la predicción de victoria local, visitante o empate en el fútbol uruguayo mediante atributos históricos causales y validación temporal. Se conserva un conjunto común de 14.705 partidos hasta 2023 y se evalúan 472 de 2024--2025. Para Naive Bayes propio y su versión de scikit-learn, la selección secuencial de suavizado y atributos produce modelos equivalentes de diez entradas, con un accuracy de 0,4788 y macro-F1 0,4454, frente a 0,4619 y 0,3564 del baseline de diez años. Por otro lado el ID3 propio alcanza 0,4661 y 0,4027, y su versión de scikit-learn sobre los mismos códigos, 0,4682 y 0,4109. El bosque aleatorio sobre tasas continuas queda en 0,4258 y 0,4219. La mayor dificultad encontrada son los empates.
\end{abstract}

\begin{IEEEkeywords}
aprendizaje supervisado, árboles de decisión, Naive Bayes, predicción deportiva, fútbol uruguayo
\end{IEEEkeywords}

% =========================================================
\section{Introducción}
La predicción de resultados del fútbol uruguayo se plantea como un problema de clasificación supervisada con tres clases, los cuales son victoria local (L), victoria visitante (V) y empate (E). A partir de registros históricos de partidos, se busca construir atributos disponibles antes del encuentro y estimar su resultado. El desafío no es únicamente ajustar un clasificador, también requiere respetar el orden temporal y evitar que el resultado que se intenta predecir intervenga en sus entradas.

Este trabajo aborda la aplicación de herramientas metodológicas a datos reales, la implementación de métodos supervisados y el análisis de sus predicciones. Se implementaron un árbol ID3 con umbral mínimo de ganancia de información, un Naive Bayes categórico con suavizado mediante tamaño equivalente de muestra y el clasificador base de proporción de victorias en los diez años anteriores. El notebook incorpora además árboles de decisión y Random Forest de scikit-learn. La comparación incluye también Naive Bayes categórico de esa biblioteca.

El procesamiento parte de seis tasas históricas de victoria y reserva 2024--2025 para evaluación. Para ambos NB, la validación seleccionó diez atributos, esas tasas más puntos y diferencia de gol por partido de cada equipo en hasta cinco antecedentes previos. Se examinan métricas por clase además de la proporción global de aciertos.

La evaluación final consolidada abarca ambos NB, el ID3 propio, dos árboles de scikit-learn y el random forest, todos sobre el mismo test. Las curvas de validación se conservan como selección de hiperparámetros y no anticipan conclusiones sobre test.

% =========================================================
% Trabajo relacionado es opcional: incorporar solo referencias verificadas
% si aportan al análisis y caben dentro de las cinco páginas totales.
\section{Conjunto de datos y preprocesamiento}
\label{sec:preprocesamiento}
% Presupuesto: aproximadamente 550--650 palabras.

\subsection{Descripción del dataset}
Se utiliza el subconjunto de fútbol uruguayo de la colección de David Schoch\cite{b1}, suministrado en un ZIP con un único CSV. El archivo contiene 15.207 filas y 17 columnas, tras el procesamiento quedan 15.177 partidos, desde el 5 de marzo de 1932 hasta el 30 de junio de 2025. Por tanto, 2025 no representa un año completo. Los campos centrales son fecha, equipos local y visitante, y goles registrados (\texttt{gh}, \texttt{ga}).

\subsection{Limpieza de datos}
La carga normaliza espacios y tipos además de exigir goles enteros no negativos. Elimina una repetición exacta. Se identifican siete grupos ambiguos de fecha y par de equipos, incluyendo localías invertidas, por lo que se excluyen sus 16 registros completos, sin elegir un marcador. Además, se descartan cinco filas de prórroga y ocho de penales. Todas estas exclusiones se aplican antes de construir historiales. Se conserva un registro de motivos y líneas del CSV. Las 175 filas admitidas con advertencias de múltiples encuentros de un equipo en una fecha mantienen sus fechas originales, no se dispone de evidencia para corregirlas. La causalidad se entiende respecto de las fechas registradas.

\subsection{Construcción de la variable objetivo}
Se adopta el resultado reglamentario sobre los registros con \texttt{full\_time} igual a \texttt{F}. La etiqueta \texttt{winner} se define como
\[
\texttt{winner} =
\begin{cases}
\text{L} & \text{si } gh>ga,\\
\text{V} & \text{si } gh<ga,\\
\text{E} & \text{si } gh=ga.
\end{cases}
\]
El diccionario de la fuente incluye prórroga y penales en los goles\cite{b1}, pero no separa marcadores reglamentarios: las ocho filas P del archivo presentan goles iguales. Se excluyen E/P por incertidumbre, sin reconstruir goles ni imponer empates, lo que limita la cobertura del estudio. Los goles actuales y \texttt{full\_time} no son entradas de los clasificadores.

\subsection{Atributos derivados}
Se construyen seis tasas: proporción de victorias en los últimos cinco partidos de cada equipo, proporción de victorias de ambos dentro del año calendario, proporción histórica de victorias del local jugando como local y proporción de victorias de ese local frente al visitante con la misma localía. Si hay menos de cinco antecedentes se usan los disponibles, la referencia anual no corresponde necesariamente a una temporada deportiva.

El cálculo recorre fechas y obtiene los atributos de todos los encuentros de un día antes de actualizar cualquier historial de ese día. Así excluye resultados simultáneos y futuros. Cuando no hay antecedentes se asigna $0{,}5$: esto evita equiparar ausencia de historial con cero victorias, pero puede coincidir con una tasa observada y no constituye una categoría independiente. Nombres de equipos y mes no integran las seis entradas; los nombres sirven para enlazar historiales. Para NB se agregan cuatro entradas: puntos por partido y diferencia de gol por partido en hasta cinco antecedentes, para local y visitante. Los puntos son 3/1/0 y la diferencia es goles a favor menos goles en contra; sin antecedentes se asignan $4/3$ y 0. Estas diez entradas, seleccionadas por validación, no incluyen goles del partido a predecir.

\subsection{Discretización}
El discretizador común produce códigos enteros para los modelos propios. Solo las dos tasas de los últimos cinco partidos (de local y de visitante) se parten con cortes fijos en $0{,}3$ y $0{,}6$, es decir, baja si $x<0{,}3$, media si $0{,}3\leq x<0{,}6$ y alta en otro caso, las restantes cuatro tasas y las cuatro entradas del NB final (puntos por partido y diferencia de gol por partido de local y visitante) usan terciles aprendidos solo con entrenamiento y reutilizados en validación o evaluación. Los cortes repetidos se eliminan, por lo que pueden resultar menos de tres intervalos; los empates entre valores tampoco garantizan grupos de igual tamaño. No existe un intervalo exclusivo para ``sin historial''. Los valores numéricos no finitos se imputan con la mediana aprendida en entrenamiento.

\subsection{División temporal de datos}
Se reservan 14.705 partidos hasta 2023 como set de entrenamiento y 472 de 2024--2025 para evaluación. Los atributos de evaluación incorporan resultados de fechas anteriores del propio período, simulando predicciones sucesivas, sin reajustar el modelo. La misma política rige para el clasificador base: todos disponen de resultados de fechas anteriores, incluidos los del bloque de evaluación.

% =========================================================
\section{Modelos}
\label{sec:modelos}
% Presupuesto: aproximadamente 400--450 palabras, más tres ecuaciones.

\subsection{Clasificador base}
TenYearWinRateClassifier compara las tasas causales de victoria de los diez años anteriores, producidas por el mismo constructor de atributos. Incluye resultados de fechas anteriores del bloque evaluado y excluye el día actual. Sin antecedentes asigna $0{,}5$, ante igualdad elige al local. Nunca predice empate y no aprende parámetros ni modifica historiales al predecir.

\subsection{Árbol de decisión propio}
La implementación utilizada es \texttt{ID3}. En cada nodo selecciona el atributo de mayor ganancia de información y crea una rama por categoría observada. Para un conjunto $S$, con proporción $p_c$ de ejemplos de clase $c$, utiliza:
\begin{equation}
H(S) = -\sum_{c \in \{L,V,E\}} p_c \log_2 p_c
\label{eq:entropia}
\end{equation}
\begin{equation}
IG(S, A) = H(S) - \sum_{v \,\in\, \operatorname{valores}(A)} \frac{|S_v|}{|S|} H(S_v)
\label{eq:ganancia}
\end{equation}
Los términos con $p_c=0$ se omiten. Cada atributo se usa como máximo una vez por rama. La recursión termina si el nodo es puro, no quedan atributos, se alcanza \texttt{max\_depth} (8 como default) o ninguna ganancia supera estrictamente \texttt{min\_info\_gain}. Este último criterio realiza prepoda. Cada nodo guarda su clase mayoritaria, que sirve de respaldo ante una categoría no observada, los empates se resuelven por orden de clase. El código contempla ramas vacías, aunque la construcción recorre únicamente categorías presentes.

\subsection{Naive Bayes propio}
MEstimateCategoricalNB estima las probabilidades de clase mediante frecuencias y supone independencia condicional entre atributos. Para el código $v$ del atributo $j$ y la clase $c$ aplica:
\begin{equation}
P(X_j=v\mid c)=\frac{n_{jvc}+m/K_j}{n_c+m}
\label{eq:m-estimate}
\end{equation}
Aquí $n_{jvc}$ cuenta coincidencias de código y clase, $n_c$ cuenta ejemplos de la clase y $K_j$ es el máximo código de entrenamiento más uno, o el mínimo declarado mediante \texttt{min\_categories} si es mayor. La grilla declara cuatro códigos en ambos NB, incluso con bines ausentes, en general $K_j$ puede variar entre atributos y un único $\alpha=m/K$ solo es equivalente con un dominio común de tamaño $K$. El cero queda reservado para desconocidos. El parámetro $m>0$ es el tamaño equivalente de muestra: reparte un peso total $m$ uniformemente entre las $K_j$ categorías, agregando $m/K_j$ a cada conteo. Se decide por $\arg\max_c[\log P(c)+\sum_j\log P(X_j=x_j\mid c)]$, con desempate E/L/V. Los códigos fuera del dominio se remiten al cero; el prior de clase $P(c)=n_c/N$ no se suaviza.

\subsection{Modelos de referencia de scikit-learn}
Se entrenan dos árboles de decisión de scikit-learn con criterio de entropía: uno recibe los mismos códigos discretos que utiliza el ID3 propio, y el otro trabaja directamente sobre las tasas continuas. El Random Forest, por su parte, combina 300 árboles entrenados sobre las tasas continuas, con semilla 42 y ponderación de clases en cada muestra bootstrap (\texttt{balanced\_subsample}).

Para comparar con el Naive Bayes propio se utiliza \texttt{CategoricalNB}, aplicado sobre la misma discretización y con cuatro códigos por atributo (incluyendo el código reservado para valores desconocidos). Este estimador calcula las probabilidades condicionales como $(n_{jvc}+\alpha)/(n_c+K_j\alpha)$. Cuando $K_j=4$ y $\alpha=m/4$, y ambos modelos comparten el mismo prior de clase empírico y las mismas entradas, esta fórmula coincide exactamente con la del NB propio (Ecuación~\ref{eq:m-estimate}). El suavizado indicado se respeta sin truncar, y el prior de clase se aprende a partir de train. En la evaluación final, ambos NB utilizan los mismos diez atributos, mientras que los árboles conservan su configuración original.

% =========================================================
\section{Experimentación}
\label{sec:experimentacion}
% PRESUPUESTO DE LA SECCIÓN COMPLETA: 180-220 palabras de TEXTO.
% Ojo: las 2 figuras de abajo son las que más espacio real consumen
% de todo el paper (aprox. 1/3 de columna cada una, ~2/3 de página en
% total entre las dos). Si el espacio aprieta más adelante, esta es
% la sección donde conviene achicar las figuras (side-by-side en una
% sola figura con 2 subplots) antes que recortar Resultados.

\subsection{Selección temporal de hiperparámetros}
Se validó en los años completos 2021, 2022 y 2023, con entrenamiento expansivo anterior a cada año. Las fechas no se dividen entre conjuntos y cada combinación ajusta un pipeline nuevo por fold. La selección utiliza macro-F1 medio con igual peso por año. El error graficado es $1-\mathrm{accuracy}$; el error de entrenamiento es de resustitución y no interviene en la elección.

Se exploraron seis umbrales ID3 entre 0 y $0{,}05$, cinco valores de $m$ entre $0{,}1$ y 1000 y cinco de $\alpha=m/4$. Ante empate exacto se favorece menor umbral o suavizado.

ID3 seleccionó $0{,}005$ (macro-F1 $0{,}3818$), ambos NB seleccionaron el menor suavizado ($m=0{,}1$, $\alpha=0{,}025$), con macro-F1 $0{,}3635$. Sus métricas coinciden, como corresponde al suavizado equivalente con cuatro códigos. Para Random Forest se conservó directamente la configuración del notebook (300 árboles sobre las tasas continuas, semilla 42, ponderación de clases en cada muestra bootstrap y hoja mínima 20), sin barrido propio de hiperparámetros. Se conservaron los dos árboles de referencia y el baseline. El notebook incluye las curvas de validación de ID3 y NB. Estos hallazgos de selección no son una evaluación independiente, el reajuste y el test de todos los modelos se presentan en la sección~\ref{sec:resultados}.

\subsection{Comparación acotada de atributos}
Con hiperparámetros fijos se comparan dos conjuntos para Naive Bayes: las seis tasas obtienen macro-F1 medio de validación $0{,}3635$ y al agregar puntos y diferencia de gol por partido, diez entradas, sube a $0{,}3858$. Se elige el mayor macro-F1 sin redondear; ante empate se prefieren seis entradas. La selección controla las columnas del discretizador y de ambos NB finales. Cada fold aprende sus cuantiles solo del train; el ID3 y el bosque mantienen seis tasas. NB conserva $m=0{,}1$ y $\alpha=0{,}025$, sin reoptimizarlos sobre diez atributos ni usar test para la selección automática. Reutilizar folds para decisiones sucesivas puede introducir optimismo.

\begin{figure}[htbp]
\IfFileExists{figuras/id3.png}{%
\centerline{\includegraphics[width=0.9\linewidth]{figuras/id3.png}}
}{\centering\fbox{\parbox{0.85\linewidth}{\centering Figura pendiente de consolidación.}}}
\caption{Error y macro-F1 de entrenamiento/validación para ID3; medias de los tres años.}
\label{fig:hiperparam-arbol}
\end{figure}

\begin{figure}[htbp]
\IfFileExists{figuras/nb_propio.png}{%
\centerline{\includegraphics[width=0.9\linewidth]{figuras/nb_propio.png}}
}{\centering\fbox{\parbox{0.85\linewidth}{\centering Figura pendiente de consolidación.}}}
\caption{Error y macro-F1 de validación para el NB propio según $m$; medias de los tres años, sobre las seis tasas iniciales.}
\label{fig:hiperparam-nb}
\end{figure}

% =========================================================
\section{Resultados}
\label{sec:resultados}
% PRESUPUESTO DE LA SECCIÓN COMPLETA: 350-400 palabras de texto,
% sumado a la tabla y a la(s) figura(s) de matriz de confusión.
% Esta es la sección con más peso relativo del informe (el enunciado
% pide explícitamente el análisis de resultados) — si sobra espacio
% de otras secciones, se lo lleva esta.

Ambos NB se ajustaron una vez con los 14.705 partidos hasta 2023 y se evaluaron sobre los mismos 472 encuentros. Los discretizadores y clasificadores permanecen fijos durante test. Las métricas agrupan todo el bloque; macro-F1 pondera E/L/V por igual y se usa cero para divisiones indefinidas.

\begin{table}[htbp]
\caption{Evaluación común sobre los 472 partidos de 2024--2025}
\label{tab:resultados}
\centering
\begin{tabular}{lrr}
\hline
Modelo & Accuracy & Macro-F1 \\
\hline
Árbol sklearn (mismos códigos) & 0,4682 & 0,4109 \\
Árbol ID3 propio & 0,4661 & 0,4027 \\
Naive Bayes propio & 0,4788 & 0,4454 \\
Naive Bayes (sklearn) & 0,4788 & 0,4454 \\
Baseline (diez años) & 0,4619 & 0,3564 \\
Random Forest & 0,4258 & 0,4219 \\
Siempre L (referencia) & 0,4025 & 0,1913 \\
Árbol sklearn (tasas crudas) & 0,3602 & 0,3514 \\
\hline
\end{tabular}
\end{table}

Los NB coinciden en las 472 predicciones; sus puntajes logarítmicos difieren como máximo $5{,}3\,10^{-15}$ y sus probabilidades, $2{,}0\,10^{-15}$. Ambos aciertan 226 partidos, ocho más que el baseline: accuracy 0,4788 frente a 0,4619 y macro-F1 0,4454 frente a 0,3564. NB obtiene el mejor resultado global de esta comparación. El árbol sklearn sobre los mismos códigos llega a 0,4682 de accuracy y 0,4109 de macro-F1; ID3 queda en 0,4661 y 0,4027, y el bosque en 0,4258 y 0,4219. Predecir siempre L da 0,4025 y 0,1913. No se evaluó significación estadística.

\begin{table}[htbp]
\caption{Métricas por clase sobre el test. Ambos NB tienen idénticos valores}
\label{tab:clases}
\centering
\begin{tabular}{llrrrr}
\hline
Modelo & Clase & Precision & Recall & F1 & $n$ \\
\hline
NB (ambos) & E & 0,3537 & 0,2214 & 0,2723 & 131 \\
 & L & 0,5475 & 0,6368 & 0,5888 & 190 \\
 & V & 0,4497 & 0,5033 & 0,4750 & 151 \\
Baseline & E & 0,0000 & 0,0000 & 0,0000 & 131 \\
 & L & 0,5021 & 0,6368 & 0,5615 & 190 \\
 & V & 0,4199 & 0,6424 & 0,5079 & 151 \\
Árbol ID3 propio & E & 0,3214 & 0,1374 & 0,1925 & 131 \\
 & L & 0,4757 & 0,7737 & 0,5892 & 190 \\
 & V & 0,5140 & 0,3642 & 0,4264 & 151 \\
Árbol sklearn (códigos) & E & 0,3281 & 0,1603 & 0,2154 & 131 \\
 & L & 0,4801 & 0,7632 & 0,5894 & 190 \\
 & V & 0,5189 & 0,3642 & 0,4280 & 151 \\
Árbol sklearn (crudas) & E & 0,2615 & 0,2595 & 0,2605 & 131 \\
 & L & 0,3990 & 0,4158 & 0,4072 & 190 \\
 & V & 0,3958 & 0,3775 & 0,3864 & 151 \\
Random Forest & E & 0,2975 & 0,3588 & 0,3253 & 131 \\
 & L & 0,5725 & 0,4158 & 0,4817 & 190 \\
 & V & 0,4261 & 0,4967 & 0,4587 & 151 \\
\hline
\end{tabular}
\end{table}

\subsection{Matrices de confusión}
Con filas reales y columnas predichas en orden E/L/V:
\[
C_{\mathrm{NB}}=
\begin{pmatrix}29&54&48\\24&121&45\\29&46&76\end{pmatrix},
\quad
C_{\mathrm{base}}=
\begin{pmatrix}0&66&65\\0&121&69\\0&54&97\end{pmatrix}.
\]
L presenta la mejor precision, recall y F1 del NB, seguida de V y de E. De sus 246 errores, 102 son empates confundidos con victorias: 54 como L y 48 como V. Las confusiones L$\leftrightarrow$V suman 91 (45 y 46), y los falsos empates, 53 (24 y 29). NB reconoce 29 de los 131 empates reales (22,1\%), frente a cero del baseline. Mantiene sus 121 aciertos en L y pierde 21 en V (76 frente a 97); el saldo son ocho aciertos adicionales. El F1 mejora en E y L, pero baja en V: la ventaja no es uniforme.

Los árboles repiten el mismo patrón con los empates:
\begin{gather*}
C_{\mathrm{ID3}}=
\begin{pmatrix}18&80&33\\24&147&19\\14&82&55\end{pmatrix}
\\[4pt]
C_{\mathrm{DT}}=
\begin{pmatrix}21&77&33\\27&145&18\\16&80&55\end{pmatrix}
\\[4pt]
C_{\mathrm{RF}}=
\begin{pmatrix}47&34&50\\60&79&51\\51&25&75\end{pmatrix}
\end{gather*}
$C_{\mathrm{ID3}}$ y $C_{\mathrm{DT}}$ aciertan 220 y 221 partidos respectivamente, con la columna L dominante (147 y 145 locales) y solo 18 y 21 empates reconocidos. El bosque reparte mejor los aciertos entre clases (79 locales, 75 visitantes y 47 empates, 201 en total sobre la diagonal), aunque su accuracy global es menor. El árbol entrenado sobre tasas continuas es el peor de los tres: acierta apenas 170 partidos y predice empate 130 veces, con toda la diagonal baja (ver Tabla~\ref{tab:resultados}).

\subsection{Análisis cualitativo}
Se ordenan los partidos por fecha, local y visitante y se toma el primero de cada par (resultado real, predicción del NB), cubriendo las nueve celdas de confusión. Es una selección reproducible e ilustrativa, no representativa; la tabla se genera en la sección~14.2 del notebook. El 11/03/2024, Montevideo Wanderers--Deportivo Maldonado terminó 0--0: NB predijo E y el baseline, L. El 17/02/2024, CA Fenix--Danubio terminó 1--2 y ambos predijeron V. Como error, el 07/04/2024 CA Cerro--Rampla Juniors Futbol Club terminó 3--0: NB predijo E y el baseline, L. Los marcadores se recuperan del ZIP limpio solo para describir los casos.

Como diagnóstico posterior se agrupan los partidos por el signo de la diferencia de puntos por partido en los últimos cinco antecedentes. Con ventaja local, NB acierta 116/216 (53,70\%); con igualdad, 15/39 (38,46\%); y con ventaja visitante, 95/217 (43,78\%). El baseline acierta 99, 21 y 98, respectivamente. La mejora se concentra en el primer grupo; en los otros dos NB obtiene menos aciertos. También cambia la composición E/L/V (62/106/48, 9/15/15 y 60/69/88), por lo que este análisis describe escenarios y no demuestra un efecto causal de la forma reciente.

Varias limitaciones ayudan a contextualizar estos resultados. La discretización pierde diferencias dentro de un mismo bin; los históricos resumen ventanas de distinta extensión y omiten sus denominadores; y una tasa de 0,5 puede ser tanto observada como neutra, por lo que no constituye evidencia de ausencia de historial. Además, tasas, puntos y goles resumen partidos superpuestos, lo que vuelve plausible cierta información compartida entre atributos, aunque no se midió su dependencia condicionada a la clase ni su efecto sobre el desempeño. Estas limitaciones no deben interpretarse como causas demostradas de los errores observados, sino como líneas abiertas para trabajo futuro.

% =========================================================
\section{Conclusiones}
El trabajo sobre Naive Bayes completó un flujo desde el ZIP original hasta limpieza, atributos causales, selección temporal, ajuste final y análisis. La selección de atributos controla el ajuste final y mantiene el mismo criterio de discretización que la búsqueda de $m$. Con dominio y suavizado equivalentes, el m-estimate propio reproduce las predicciones de \texttt{CategoricalNB}. NB mejora sobre el baseline en accuracy (0,4788 frente a 0,4619) y macro-F1 (0,4454 frente a 0,3564), aunque no en todas las clases. El empate sigue siendo la principal dificultad.

La evaluación es secuencial y 2025 está incompleto: no representa un pronóstico simultáneo de dos años ni evidencia causal. La selección automática usa validación temporal, pero el test de 2024--2025 ya fue inspeccionado en versiones anteriores; esta corrección de consistencia no constituye una evaluación confirmatoria sobre un holdout intacto. En estos 472 partidos, NB obtiene el mayor accuracy y macro-F1, por delante de los árboles de referencia. En la comparación concreta de árboles, discretizar las seis tasas mejora accuracy frente a los valores continuos (0,3602); este resultado no establece un efecto general. Ningún modelo resuelve satisfactoriamente el empate.

Estas cifras reflejan, más que una limitación puntual de NB o ID3, una dificultad estructural del problema. Un resultado de fútbol depende de factores que ninguna tasa histórica captura --lesiones, decisiones arbitrales, clima, estado anímico del plantel-- y los dos encaran el problema de forma distinta: el NB supone, de manera explícita, independencia condicional entre atributos dada la clase, mientras que el árbol sí puede capturar interacciones mediante divisiones sucesivas, aunque solo sobre particiones alineadas con las categorías discretas de cada atributo (expresar que una racha positiva del local importe menos si el rival tiene mejor historial head-to-head exige dividir ambos atributos en el orden y el corte adecuados). La discretización en pocos bines agrava esto, ya que dos partidos con tasas apenas distintas (0,45 y 0,55) pueden terminar en la misma categoría o en categorías opuestas según dónde caiga el corte. A esto se suma que el empate, por definición, es un resultado de "equilibrio" entre dos equipos, y las tasas de victoria usadas como atributos no fueron diseñadas para capturar paridad sino para capturar dominancia; es esperable, entonces, que sea la clase peor predicha en cualquier modelo que dependa únicamente de ese tipo de señal.

\section*{Uso de Inteligencia Artificial}
En la parte de NB se utilizó Codex para asistir en implementación, depuración, pruebas, análisis de resultados y revisión de la redacción; también se empleó en esta corrección de consistencia y su verificación. Para ID3 se utilizó Claude. Los autores mantienen la responsabilidad sobre la comprensión y validación del trabajo.

\begin{thebibliography}{00}
\bibitem{b1} D. Schoch, ``football-data,'' GitHub repository. [Online]. Available: https://github.com/schochastics/football-data
% TODO: agregar referencias adicionales si se citan en Trabajo relacionado.
\end{thebibliography}

\end{document}
